Suppose $(M, g)$ and $(N,h)$ are piecewise smooth Riemannian manifolds.   If $f: M \rightarrow N$ is piecewise smooth,  we define the $k$-dilation of $f$ as the infimal $\Lambda$ so that for every $k$-dimensional submanifold $\Sigma \subset M$, 

$$ \Vol_k( f(\Sigma) ) \le \Lambda \Vol_k (\Sigma). $$

\noindent We write the $k$-dilation of $f$ as $\Dil_k(f)$.   Equivalently,  $\Dil_k(f) = \sup_{x \in M} \| \Lambda^k df \|$.

Prove that there is a constant $k > 0$ so that the following holds.

Suppose that $R \subset \RR^4$ is a 4-dimensional rectangle with side lengths $R_1 \le R_2 \le R_3 \le R_4$ and $S$ is a 4-dimensional rectangle with side lengths $S_1 \le S_2 \le S_3 \le S_4$.   Suppose that $f: (R, \partial R) \rightarrow (S, \partial S)$ is a piecewise smooth map with degree 1 and with $\Dil_2(f) \le 1$.   Suppose that $R_1 \le k S_1$.   Then either

$$ R_3 R_4 > k S_3 S_4,  or $$

$$ R_1 R_2 R_3 R_4 > k S_1 S_2^{1/2} S_3^{3/2} S_4. $$